\section*{2. Personas}
\addcontentsline{toc}{section}{2. Personas}

\subsection*{Definition}
\addcontentsline{toc}{subsection}{Definition}
A user persona imagines a specific type of user for a given platform. e.g. SETUP, a student services app, may consider students, lecturers, and/or guests. Product expectations may vary depending on user background, education and skillset.\\ 

By conceptualising the types of users this product seeks to attract; personas assess what features are best to proritise. It also suggests what features may require further development, or scrapping entirely, in order for the product to be suitable for its demographics.

\subsection*{1. PLACEHOLDER, a COURSENAME student}
\addcontentsline{toc}{subsection}{Persona 1}
incl. relevant job/course background details
persona goals/ambitions\\
educational/technical skillset; what would make ui intuitive?\\
relevance of persona to SETUP app
 
\subsection*{2. PLACEHOLDER, a COURSENAME lecturer}
\addcontentsline{toc}{subsection}{Persona 2}
incl. relevant job/course background details
persona goals/ambitions\\
educational/technical skillset; what would make ui intuitive?\\
relevance of persona to SETUP app

\subsection*{3. Maria Angelo, Parent, Nursing Student}
\addcontentsline{toc}{subsection}{Persona 3}

Maria Angelo, 46 years old, mother of two daughters. She is a nurse who works in a local hospital, aiding patients in need of care. She graduated out of nursing school when she was 22 years old in the Philippines. By 30 years old she got the opportunity to work
in Ireland, where she met her lovely husband 2 years after. They live in the suburb with their daughters, a 12 years old and and a 14 year old. One who is starting first year of secondary school while the other is in third year. Her husband, who is an accountant, works in a business company
and is rarely around the house as he's always busy with his 9 to 5 job. Maria has been in the nursing field for 23 years and she wants to specialize in a different area where she needs a certification.\\ 

She decides to go back to college and start studying again. Maria is a busy woman, now that she's back in education, works in the hospital and sometimes would go overtime. She packs lunch for her daughters and drops them off to school. Her husband would alternate around giving their daughters lifts. 
After she finishes college, she picks up her daughters up and would go straight to work in the evenings. Maria loves keeping herself busy but from time to time she would forget what to do next. She would miss deadlines for her assignments and would have to contact her lecturers to have an extension. She doesn't make schedules, she would rather write the information on the palm of her hand. When Maria realised that she developed additional stress into her life, she is 
willing to start being organise.\\

With the help of SETUP, she can create her schedules, check lists of her tasks and additional information. This will assist her time schedule, as it includes an automated timetable specifically for the user. There's a board for the user to write up their personal to-do list and it traces deadlines that is set by Moodle. She can email her lecturers through the app as well as receiving them. 
Instead of going through different applications for peculiar tasks, she can do it all in one application.\\
 
The app is set to handle these jobs. The built in integrations and neat GUI will help her stay on top and well arrange to set personal schedules, tick off her to-do lists and communicate with her lecturers with deadlines. 


\subsection*{4. PLACEHOLDER, a COURSENAME student with accessibility concerns}
\addcontentsline{toc}{subsection}{Persona 4}
incl. relevant job/course background details
persona goals/ambitions\\
educational/technical skillset; what would make ui intuitive?\\
relevance of persona to SETUP app.